\subsubsection{Modèle des bandes : un outils pour comprendre l'optimisation des propriétés rédox}

%La structure de bande nous donne des informations sur la stabilité électronique d'un édifice moléculaire, l'origine moléculaire la tension à courant nul (aussi appelée OCV ou OCP, soit open-circuit voltage (or potential) $V_{oc}$), les atomes impliqués dans les processus d'oxydo-réduction, $\ldots$

\begin{comment}
\begin{center}
\begin{minipage}[]{0.2\textwidth}
\includegraphics[height=3cm]{Ex17/Fig01}
\end{minipage}\begin{minipage}[]{0.3\textwidth}
\captionof{figure}{Structure de bande de différentes cathodes }
\label{fig:Ex16-Fig02}
\end{minipage}
\end{center}

\Q Rappelez quelles sont les orbitales atomiques impliquées dans les bandes d'énergie étroites associés à chacune des cathodes.

\Q Justifier en quoi l'existente d'une bande d'énergie étroite est une garantie de stabilité des cathodes.

\Q Pour chacune des cathodes, calculer le nombre d'électron de valence impliqué dans l'orbitale de plus basse énergie.

\paragraph*{Relation entre la tension à courant nul et le diagramme de bande}

Pour les 4 matériaux présentés ci-dessus, l'OCV a été mesurées pour différents degrés de lithiation.

\begin{center}
\begin{minipage}[]{0.2\textwidth}
\includegraphics[height=3cm]{Ex17/Fig02}
\end{minipage}\begin{minipage}[]{0.3\textwidth}
\captionof{figure}{Profil de potentiel en décharge, par rapport au couple \ce{Li+}/\ce{Li^0} , des cathodes \ce{Li_x C_6} , \ce{Li_x Ti S_2} and \ce{Li_x [Ti_2] S_4} , \ce{Li_x Co O_2} , and \ce{Li_x Co PO_4}.}
\label{fig:Ex16-Fig02}
\end{minipage}
\end{center}


\Q \`{A} partir du diagramme de potentiel, explicitez l'origine électronique de cette variation de tension.

\paragraph*{Substitution anionique}

\Q

\paragraph*{Substitution cationique}

\Q

\paragraph*{Bilan}

\Q En quelques mots, résumez les outils dont dispose le chimiste pour élaborer des matériaux d'électrodes efficace.

\end{comment}
